\begin{textsample}{POS Dim 4 – 2011\_c6 – Score 2.00 – 2011\_c6/enade\_2011\_geo\_ada47c37e1.txt}  \label{ex:f4_pos_300}
CAPÍTULO 3 ANÁLISE TÉCNICA DA PROVA Este capítulo tem por objetivo apresentar o desempenho dos estudantes concluintes de Geografia ( Bacharelado e Licenciatura ) no ENADE/2011.
Para isso, foram calculadas as estatísticas básicas da prova em seu todo, bem como as estatísticas dos componentes relacionadas à Formação Geral, ao de Conhecimento Específico da Área e das questões discursivas isoladamente.
Além disso, as análises das questões objetivas do componente de Conhecimento Específico são apresentadas separadamente para concluintes do Bacharelado e da Licenciatura, já que as questões propostas para cada grupo foram diferentes.
Nas tabelas, são apresentados o tamanho da população inscrita e de presentes, e as seguintes estatísticas das notas⁸ : média do desempenho na prova, erro padrão da média, desvio padrão, nota mínima, mediana e nota máxima.
As estatísticas apresentadas neste capítulo contemplam o total de estudantes concluintes da área de Geografia ( Bacharelado e Licenciatura ) em 2011 do Brasil e, separadamente, por Grande Região.
Foram calculadas tendo -se em vista as seguintes agregações : ( a ) as Grandes Regiões e o país como um todo ; ( b ) a Categoria Administrativa ; e ( c ) a Organização Acadêmica.
Em relação aos gráficos de distribuição de notas, o intervalo considerado foi de 10 unidades, aberto à esquerda e fechado à direita, com exceção do primeiro intervalo, [ 0 ; 10 ], fechado em ambos os extremos.
Para os gráficos de distribuição das notas das questões discursivas, foram consideradas mais duas categorias : questão em branco e nota zero.
Todos os gráficos de distribuição de notas permitem a comparação dos resultados por habilitação : Bacharelado e Licenciatura. 3.1 ESTATÍSTICAS BÁSICAS DA PROVA 3.1.1 Estatísticas Básicas Gerais A Tabela 3.1 apresenta as estatísticas básicas da prova por grande Região dos estudantes concluintes de Geografia ( Bacharelado e Licenciatura ).
A população total de inscritos foi de 13.386.
Destes, 9.942 estiveram presentes, sendo 25,7 \% o índice de não comparecimento.
A Região de maior abstenção foi a Norte ( 30,5 \% ), e a de menor abstenção foi a Centro-Oeste ( 21,9 \% ). ⁸ Essas estatísticas e outras estão definidas no Capítulo 1.
A média das notas da prova como um todo ( nas seções seguintes serão analisados os componentes de Formação Geral e de Conhecimento Específico ) foi 39,4, sendo que os alunos da região Sudeste obtiveram a média mais baixa ( 37,8 ), e os da região Sul obtiveram a média mais alta ( 41,1 ).
As demais médias foram : 38,9 na região Norte ; 40,7 na região Nordeste ; e 39,5 na região Centro-Oeste.
O desvio padrão para o Brasil como um todo foi 17,8, sendo o maior desvio padrão encontrado na região Sudeste ( 20,4 ), e o menor, na região Centro-Oeste ( 14,0 ), indicando uma dispersão um pouco menor das notas desta última região.
A região que obteve a maior nota máxima foi a Nordeste ( 86,3 ), ao passo que a região que atingiu a menor nota máxima foi a Centro-Oeste ( 76,1 ).
A mediana do Brasil como um todo foi 41,9, sendo a maior mediana obtida na região Sul ( 43,6 ) e a menor obtida na Norte ( 40,1 ).
A nota mínima foi zero em todas as regiões, sem exceção.
O comportamento das notas dos estudantes de todo o Brasil pode ser observado no Gráfico 3.1 que apresenta a distribuição das mesmas em colunas diferentes para alunos de Bacharelado e de Licenciatura.
As distribuições são unimodais, com um segundo máximo local no primeiro intervalo [ 0 ; 10 ].
O intervalo modal principal dos concluintes de Bacharelado é ( 50 ; 60 ] e dos de Licenciatura ( 40 ; 50 ].
Os coeficientes de assimetria das distribuições das notas das duas habilitações são negativos : – 0,58 para o Bacharelado e – 0,64 para a Licenciatura.
As distribuições por Grande Região também apresentam assimetria negativa, concentração pouco maior do lado direito do histograma e mais espalhada do lado esquerdo.
Os Gráficos 3.2, 3.3 e 3.4 apresentam informações referentes à média da nota final dos Participantes, desagregando os resultados de acordo com, respectivamente, as Grandes Regiões do país, a Categoria Administrativa e a Organização Acadêmica.
Os gráficos apresentam o valor da média das notas como uma barra e os extremos do intervalo de confiança de 95 \% como linhas verticais unidas por uma linha horizontal na forma da letra H maiúscula.
Considerando -se o gráfico de notas segundo Grande Região ( Gráfico 3.2 ), observa -se que existe diferença estatisticamente significativa ao nível de 95 \% entre a maior média, obtida na região Sul ( 41,1 ) e a menor, obtida na região Sudeste ( 37,8 ).
Levando -se em conta os estudantes por Categorias Administrativas ( Gráfico 3.3 ), observa -se que existe diferença estatisticamente significante entre as médias das notas das IES Públicas e Privadas.
Pode ser observado que a média dos alunos de IES Públicas ( 38,6 ) é menor do que a dos alunos de IES Privadas ( 42,0 ).
Tendo como foco o Gráfico 3.4, que apresenta as notas médias das provas segundo Organização Acadêmica, constata -se que existe diferença estatisticamente significativa ao nível de 95 \% nas médias das notas dos estudantes provenientes de Centros Universitários ( 42,0 ) em relação aos de Universidades ( 39,2 ) e Faculdades ( 39,9 ). 3.1.2 Estatísticas Básicas no Componente de Formação Geral A Tabela 3.2 apresenta as estatísticas básicas em relação ao componente da prova que avalia a Formação Geral dos estudantes concluintes.
Os alunos de todo Brasil obtiveram desempenho médio de 47,0.
Quanto à variabilidade, o desvio padrão das notas dos estudantes do Brasil como um todo foi 23,0.
A maior média foi obtida na região Centro-Oeste ( 49,4 ), e a menor, na região Sudeste ( 43,8 ).
As demais médias foram : 48,1 na região Norte ; 48,9 na região Nordeste ; e 48,5 na região Sul.
Já o maior desvio padrão foi obtido na região Sudeste ( 25,6 ) e o menor na região Centro-Oeste ( 19,2 ).
Os demais desvios padrões foram : 20,0 na região Norte ; 21,5 na região Nordeste ; e 22,9 na região Sul.
A maior nota no componente de Formação Geral da prova do ENADE foi obtida por pelo menos um aluno das regiões Norte e Nordeste ( 100,0 ) enquanto que a menor nota máxima foi obtida na região Centro-oeste ( 94,0 ).
Nas outras regiões as notas máximas foram : 98,0 na região Sudeste e 99,0 na região Sul.
A mediana do Brasil como um todo foi 50,5, sendo a menor mediana encontrada na região Sudeste ( 49,0 ) e a maior encontrada nas regiões Nordeste e Sul ( 52,5 ).
A nota mínima nesta parte foi zero em todas as regiões, sem exceção.
O Gráfico 3.5 propicia a avaliação do desempenho dos estudantes no componente de Formação Geral a partir do histograma da distribuição das notas correspondentes ao Bacharelado e à Licenciatura.
As distribuições são bimodais : a moda principal da distribuição de notas dos alunos da Licenciatura está no intervalo ( 50 ; 60 ] e a moda secundária no primeiro intervalo, [ 0 ; 10 ].
A situação se inverte para a distribuição dos alunos de Bacharelado, a moda principal está no intervalo [ 0 ; 10 ] e a secundária, mais acima, no intervalo ( 60 ; 70 ].
Nota -se, que no Gráfico 3.5 as notas apresentam uma maior dispersão do que no Gráfico 3.1 ( distribuição das notas da prova ), confirmado pela comparação dos desvios padrões : 17,8 para a nota da prova como um todo e 23,0 para o componente de Formação Geral.
Para o componente de Formação Geral, o coeficiente de assimetria da distribuição das notas dos estudantes, como na prova como um todo, também é negativo ( – 0,55 ).
Desagregando por habilitação observa -se que o coeficiente da distribuição dos concluintes do Bacharelado tem menor valor absoluto ( – 0,36 ) do que o da Licenciatura ( – 0,59 ), o que indica uma assimetria à esquerda mais acentuada para o segundo grupo.
Em todas as Grandes Regiões os histogramas também possuem assimetria negativa, tomando -se os alunos em conjunto ou por habilitação.
Nos Gráficos 3.6, 3.7 e 3.8 são apresentadas as informações referentes ao desempenho dos Concluintes no componente de Formação Geral, em diferentes agregações : Grande Região do país, Categoria Administrativa e Organização Acadêmica.
Observa -se pelo Gráfico 3.6 que existe diferença estatisticamente significativa apenas entre a menor média das notas no Componente de Formação Geral, obtida na região Sudeste ( 43,8 ), e as médias das demais Grandes Regiões do país.
No Gráfico 3.7, que representa as notas médias no Componente de Formação Geral segundo Categoria Administrativa, observa -se que existe diferença estatisticamente significativa entre as médias.
Para o componente de Formação Geral, da mesma forma que para a prova como um todo, os concluintes de Geografia ( Bacharelado e Licenciatura considerados em conjunto ) das IES Públicas ( 45,8 ) obtiveram uma média menor do que os das IES Privadas ( 50,8 ).
Considerando -se o tipo de Organização Acadêmica, nota -se, no Gráfico 3.8, uma diferença estatisticamente significativa entre a menor média e as demais.
Nas Universidades ( 46,4 ) a média é mais baixa do que em Centros Universitários ( 50,5 ) e Faculdades ( 50,0 ). 3.1.3 Estatísticas Básicas do Componente de Conhecimento Específico A Tabela 3.3 apresenta as estatísticas básicas referentes ao componente de Conhecimento Específico da área de Geografia ( Bacharelado e Licenciatura ).
A média do desempenho dos alunos do Brasil como um todo foi 36,9.
A maior média foi obtida na região Sul ( 38,6 ), e a menor, nas regiões Norte e Sudeste ( 35,8 ).
As demais médias foram : 37,9 na região Nordeste e 36,2 na região Centro-Oeste.
Quanto à variabilidade das notas, o desvio padrão do Brasil como um todo foi 17,6, sendo o maior desvio padrão observado na região Sudeste ( 20,0 ) e o menor na região Norte ( 14,2 ).
Os demais desvios foram : 16,3 da região Nordeste ; 18,1 da região Sul ; e 14,3 da região Centro-Oeste.
A mediana das notas dos estudantes de todo o Brasil foi 39,3.
A maior mediana ocorreu na região Sul ( 40,9 ) e a menor na região Norte ( 36,0 ).
As demais medianas foram : 39,3 nas regiões Nordeste e Sudeste ; e 36,2 na região Centro-Oeste.
A nota máxima do Brasil como um todo foi 90,3, obtida por pelo menos um aluno da região Sul.
As demais notas máximas foram : 72,0 na região Norte ; 89,0 na região Nordeste ; 84,9 na região Sudeste ; e 72,2 na região Centro-Oeste.
A nota mínima foi zero em todas as regiões.
Assim como os Gráficos 3.1 e 3.5, o Gráfico 3.9, apresentado a seguir, proporciona uma avaliação do desempenho de concluintes dos cursos de Bacharelado e de Licenciatura em relação ao componente de Conhecimento Específico com um histograma da distribuição das notas correspondentes.
Estas também são distribuições bimodais e o grupo modal principal das notas da Licenciatura é o intervalo ( 30 ; 40 ] e das notas do Bacharelado é ( 50 ; 60 ].
O máximo local, a moda secundária, é a mesma para as duas habilitações, o primeiro grupo de notas, [ 0 ; 10 ].
O coeficiente de assimetria da distribuição das notas do componente de Conhecimento Específico é negativo tanto pra a distribuição das notas do curso de Bacharelado ( – 0,49 ) quanto para as de Licenciatura ( – 0,38 ).
Nota -se pelo histograma ( Gráfico 3.9 ) que para ambas as habilitações a cauda é um pouco mais pesada à esquerda.
Para todas as regiões os coeficientes de assimetria são negativos para os dois cursos.
A assimetria à esquerda mais acentuada é a da distribuição de notas dos alunos de Bacharelado da região Nordeste ( coeficiente – 0,70 ) e a mais leve encontrada na região Centro-Oeste ( – 0,21 ), para alunos de Licenciatura.
Os Gráficos 3.10, 3.11 e 3.12 apresentam uma comparação dos resultados em relação à Grande Região do país, à Categoria Administrativa e à Organização Acadêmica, agora levando em conta o desempenho de alunos no componente de Conhecimento Específico da prova.
Pelo Gráfico 3.10, observa -se que existe diferença estatisticamente significativa da maior média das notas no componente de Conhecimento Específico, da região Sul ( 38,6 ), em relação às médias das regiões Centro-Oeste ( 36,2 ), Sudeste e Norte ( 35,8 ).
As médias das regiões Norte e Sudeste ( 35,8 ) é significativamente menor do que todas as outras, com exceção da região Centro-Oeste.
Quanto à Categoria Administrativa ( gráfico 3.11 ), observa -se um comportamento semelhante àquele da parte de Formação Geral e à prova como um todo, ou seja, existe diferença estatisticamente significativa entre as médias das IES Públicas ( 36,2 ) e Privadas ( 39,0 ), sendo a maior média obtida por alunos de IES Privadas.
Quanto ao Gráfico 3.12, observa -se que existe diferença estatisticamente significativa, ao nível de 95 \%, entre as notas dos concluintes de Centros Universitários ( 39,1 ), maior média, em relação aos das demais organizações acadêmicas.
Entre as médias dos que fizeram o curso em Universidades ( 36,8 ) e Faculdades ( 36,5 ) a diferença não é significativa. 3.2 ANÁLISE DAS QUESTÕES OBJETIVAS 3.2.1 Componente de Formação Geral A Tabela 3.4 apresenta as estatísticas básicas relativas às oito questões objetivas do componente da prova que abrange a Formação Geral dos estudantes.
A média do Brasil foi 48,2.
A menor média foi encontrada na região Sudeste ( 45,4 ) e a maior na região Nordeste ( 50,1 ).
As demais médias foram 49,1 na região Norte e 49,6 nas regiões Sul e Centro-Oeste.
O desvio padrão do Brasil foi 24,3, sendo o maior desvio padrão encontrado na região Sudeste ( 26,8 ) e o menor na região Centro-Oeste ( 21,0 ).
Os demais desvios foram : 21,9 na região Norte ; 22,6 na região Nordeste ; e 24,5 na região Sul.
As medianas ( 50,0 ), as notas máximas ( 100,0 ) e as notas mínimas ( 0,0 ) foram iguais para todas as regiões.
A Tabela 3.5 apresenta os índices de facilidade e discriminação ( ponto bisserial ) das questões objetivas do componente de Formação Geral.
Quanto ao índice de facilidade, foram usadas as seguintes cores para diferenciar o nível de dificuldade da questão :  Azul para as questões classificadas com índice muito fácil ( índice \&gt ; =0,86 ), verde para as questões classificadas com índice fácil ( 0,61 a 0,85 ), amarelo para as questões classificadas com médio ( 0,41 a 0,60 ), vermelho para as questões classificadas com difícil ( 0,16 a 0,40 ) e roxo para as questões classificadas com muito difícil ( \&lt ; =0,15 ).
Já quanto ao índice de discriminação, foram usadas as seguintes cores para qualificar a questão :  As questões classificadas com índice fraco receberam a cor vermelho ( índice \&lt ; =0,19 ), as classificadas com médio receberam a cor amarelo ( 0,20 a 0,29 ), as classificadas com bom receberam a cor verde ( 0,30 a 0,39 ) e as classificadas com muito bom ( \&gt ; =0,40 ) receberam a cor azul.
As questões objetivas do componente de Formação Geral, segundo o índice de facilidade, foram assim avaliadas : das oito questões, nenhuma teve o índice de facilidade classificado como muito fácil.
Três questões foram tidas como fáceis, por terem índice de acertos situado na faixa entre 0,61 e 0,85 ( de 61,0 \% a 85,0 \% de acertos ).
Duas questões foram consideradas de dificuldade média, situando -se no intervalo entre 0,41 e 0,60 do índice de facilidade, ou seja, houve entre 41,0 \% e 60,0 \% de acertos.
Outras três questões foram classificadas na categoria difícil, situando -se no intervalo entre 0,16 e 0,40.
Por fim, nenhuma das questões apresentou menos de 15 \% de acertos, razão pela qual seriam classificadas como muito difíceis.
Como já comentado, para análise das questões objetivas relativas à Formação Geral segundo o poder de discriminação, utilizou -se, o índice de discriminação ponto bisserial.
Nesta análise as questões foram assim avaliadas : sete das oito questões apresentaram índices acima ou igual a 0,40 e, assim, foram classificadas com índice muito bom para esse grupo de estudantes.
Uma questão teve bom índice de discriminação, com valor entre 0,30 e 0,39 sendo classificada com índice bom para esse grupo de estudantes.
Nenhuma questão teve nível fraco de discriminação.
O índice de facilidade variou de 0,18 a 0,68, e o de discriminação, de 0,33 a 0,65.
As sete questões com índices de discriminação muito bom, as de números 1 a 7, figuraram entre os diversos níveis de dificuldade desse conjunto : três classificadas na categoria fácil ( questões 1, 3 e 6 ) do índice de facilidade, duas na categoria médio ( questões 4 e 5 ) e duas na categoria difícil ( questões 2 e 7 ).
Em particular, a questão 3 foi a que apresentou o maior poder discriminatório, com índice 0,65, e foi também uma das mais fáceis, com uma proporção de 0,62 acertos.
O máximo de acertos foi alcançado pela questão 6 com um índice de facilidade de 0,68.
A questão de número 4 apresentou um índice de facilidade de 0,47, ou seja, quase metade dos estudantes conseguiu resolvê -la, dentro do universo de participantes.
Seu índice de discriminação foi muito bom ( 0,44 ).
Já a questão 8 obteve índice de discriminação bom, 0,33, e seu índice de facilidade foi difícil ( 0,18 ).
O Gráfico 3.13, para exemplificar, analisa o comportamento da questão de número 3 de Formação Geral.
Trata -se da segunda questão mais fácil e a que obteve o maior índice de discriminação dessa parte da prova.
Neste gráfico, cada uma das cinco curvas representa o percentual de respostas em determinada alternativa da questão, em função do número de acertos dos estudantes nesta parte da prova ( Formação Geral/Múltipla Escolha ), antes de possíveis eliminações pelo critério do ponto bisserial.
Em princípio, a soma das escolhas possíveis deveria ser igual a 100 \%.
Não é este o caso, pois não aparecem no gráfico as questões deixadas em branco ou com múltiplas respostas.
Como foram oito as questões, os valores variam de 0 a 8 acertos.
A curva em vermelho corresponde à alternativa E, a correta para esta questão.
Assim, observa -se que entre os estudantes com menor número de acertos, nessa parte do exame, a situação mais frequente foi a escolha de uma das alternativas incorretas : a alternativa A ( em azul ).
Na medida em que o número de acertos aumenta, indicando desempenho melhor nesta parte da prova, aumenta concomitantemente a proporção de estudantes que selecionaram a alternativa correta E, atingindo 100 \% para os estudantes com 8 acertos.
Essa análise permite verificar como a questão discriminou os grupos de desempenho, justificando o alto índice obtido na questão.
Os gráficos relativos às demais questões de Formação Geral constam do Anexo I. 3.2.2 Componente de Conhecimento Específico – Geografia Bacharelado A Tabela 3.6 apresenta as estatísticas básicas em relação às questões objetivas do componente de Conhecimento Específico da prova de Geografia – Bacharelado, por Grande Região.
A média do Brasil deste componente foi de 42,4.
A menor média foi observada na região Sudeste ( 39,5 ) e a maior na região Nordeste ( 48,3 ).
O desvio padrão de todo o Brasil foi 23,6, sendo o menor desvio padrão encontrado na região Centro-Oeste ( 18,1 ) e o maior na região Sudeste ( 25,2 ).
A mediana de todo o Brasil foi 48,1, a mesma encontrada na região Sudeste.
A região Norte apresentou mediana com valor menor ( 40,7 ) e para as demais regiões, Nordeste, Sul e Centro-Oeste, o valor da mediana foi maior ( 51,9 ).
A nota máxima da prova foi 92,6, obtida nas questões objetivas do componente de Conhecimento Específico, por pelo menos um aluno da região Sudeste.
A menor nota máxima foi 77,8 na região Norte.
A Tabela 3.7 apresenta os índices de facilidade e discriminação ( ponto bisserial ) das questões objetivas do componente de Conhecimento Específico para os estudantes de Geografia – Bacharelado.
Para facilitar a diferenciação das questões usou -se as mesmas cores das Tabelas 3.5 para as diferentes classificações dos índices de facilidade e de discriminação.
Dentre as questões objetivas da parte da prova relativa ao componente de Conhecimento Específico do Bacharelado, nenhuma foi anulada pela Comissão.
Desse modo, a classificação quanto ao índice de facilidade foi estabelecida com base em todas as 27 questões : de 9 a 25 e de 36 a 45.
A partir dos índices obtidos, nenhuma foi considerada muito difícil e também não houve questão classificada como muito fácil.
Menos da metade das questões objetivas da prova foi considerada difícil : 11 das 27 questões.
A maioria das demais questões ( 15 ) foram classificadas como de dificuldade média, na faixa de 0,41 a 0,60 do índice de facilidade, e apenas uma foi considerada fácil, índice entre 0,61 e 0,85.
Já quanto aos índices de discriminação das questões objetivas do componente de Conhecimento Específico da prova, tem -se como resultado a seguinte classificação : seis das 27 questões foram consideradas como boas, enquanto 20 delas tiveram índice de discriminação muito bom.
Assim, para a maioria das questões – 26 em 27 – os índices de discriminação foram bons ou muito bons.
Apenas uma, questão 36, foi classificada como média.
Constata -se, assim, que a prova – no que se refere ao componente de Conhecimento Específico – possuía capacidade muito boa de discriminar entre aqueles que dominam ou não o conteúdo.
Dentre as 20 questões que alcançaram os maiores índices de discriminação, muito bom, situando -se no intervalo de 0,40 a 0,50 do índice, apenas uma ( questão 45 ) foi classificada na categoria fácil, quanto ao índice de facilidade.
As demais se distribuíram nas categorias médio ( questões 9, 10, 11, 13, 16, 17, 21, 23, 24, 25, 37, 39, 40 e 43 ) e difícil ( questões 14, 18, 19, 42 e 44 ) do índice de facilidade.
A questão de número 22 foi a mais difícil dentre as 27 questões específicas, com baixo índice de facilidade, apenas 23,0 \% de acertos.
Apesar disso, essa questão apresentou poder discriminatório bom, 0,31.
Destaca -se, também, a questão 36, com o menor índice de discriminação, 0,24, e que também teve baixo índice de facilidade, 0,26, o que, em termos percentuais, corresponde a 26,0 \% de estudantes que responderam acertadamente.
Assim, como nenhuma questão teve índice de discriminação abaixo de médio, não houve questão desconsiderada no cômputo da nota final.
A título de exemplo das análises do comportamento das questões objetivas, o Gráfico 3.14 analisa a questão 45 do componente de Conhecimento Específico.
Esta foi considerada pelos estudantes avaliados a questão mais fácil da prova, apresentando índice de facilidade 0,70, ou seja, 70,0 \% dos estudantes assinalaram acertadamente a opção A, correspondente ao gabarito.
Seu índice de discriminação foi igual a 0,78, classificado como muito bom, também sendo esta questão que apresentou o maior índice discriminatório.
Neste gráfico, cada uma das cinco curvas representa o percentual de respostas em determinada alternativa da questão 45, em função do número de acerto dos estudantes nesta parte da prova, antes de possíveis eliminações de questões pelo critério do ponto bisserial.
A alternativa correta A, representada no gráfico pela curva em azul, foi escolhida em maiores proporções pelos alunos com desempenho melhor nesta parte da prova.
Já as alternativas incorretas, também denominadas distratores, foram selecionadas principalmente por aqueles com notas mais baixas.
Neste caso também a soma não é sempre 100 \% por causa das questões não respondidas ou com mais de uma opção marcada.
Aqueles com nota zero, na sua quase totalidade deixaram esta questão em branco ou marcaram mais de uma alternativa, comportamento considerado inválido.
A proporção de alunos que selecionou a resposta correta A aumenta gradativamente, chegando a atingir 100 \% para mais de 20 acertos, enquanto a proporção dos que escolheram alternativas incorretas decai, a partir de seis acertos, como função do número de acertos nesta parte da prova.
Os gráficos relativos às demais questões do componente de Conhecimento Específico constam do Anexo I. 3.2.3 Componente de Conhecimento Específico – Geografia Licenciatura A Tabela 3.8 apresenta as estatísticas básicas em relação às questões objetivas do componente de Conhecimento Específico da prova de Geografia - Licenciatura, por Grande Região.
A média do Brasil deste componente foi de 41,1.
A menor média foi observada na região Centro-Oeste ( 39,5 ) e a maior na região Sul ( 43,3 ).
O desvio padrão de todo o Brasil foi 18,2, sendo o maior desvio padrão encontrado na região Sudeste ( 20,6 ) e o menor na região Centro-Oeste ( 15,2 ).
A mediana de todo o Brasil foi 42,3, também encontrada para as regiões Nordeste e Sudeste.
A maior mediana ocorreu na região Sul ( 46,2 ) e a menor na região Centro-Oeste ( 38,5 ).
A nota máxima da prova foi 96,2, obtida por pelo menos um concluinte de Licenciatura da região Sul, nas questões objetivas do componente de Conhecimento Específico.
A menor nota máxima ( 80,8 ) ocorreu na região Centro-Oeste.
A Tabela 3.9 apresenta os índices de facilidade e discriminação ( ponto bisserial ) das questões objetivas do componente de Conhecimento Específico para a prova de Geografia – Licenciatura ( questões 9 a 35 ).
Para facilitar a diferenciação das questões usou -se as mesmas cores das Tabelas 3.5 e 3.7 para as diferentes classificações dos índices de facilidade e de discriminação.
Dentre as questões objetivas da parte da prova relativa ao componente de Conhecimento Específico de Geografia - Licenciatura, nenhuma foi anulada pela Comissão.
Desse modo, a classificação quanto ao índice de facilidade foi estabelecida com base em todas as 27 questões.
A partir dos índices obtidos, observa -se que 12 das 27 questões foram consideradas difíceis e que nenhuma foi classificada como muito difícil.
Também não houve questão classificada como muito fácil, ao passo que duas foram tidas como fáceis, na faixa de 0,61 a 0,85 do índice de facilidade, e as outras 13 consideradas médias, entre 0,41 e 0,60.
Já quanto aos índices de discriminação das questões objetivas do componente de Conhecimento Específico da prova, tem -se como resultado a seguinte classificação : 11 das 27 questões foram consideradas como muito boas, enquanto 12 delas tiveram índice de discriminação bom.
Assim, para a maioria das questões – 23 em 27 – os índices de discriminação foram bons ou muito bons.
Dentre as demais, três delas foram classificadas como médias e apenas uma como fraca, sendo quatro, por conseguinte, a quantidade de questões nos dois patamares mais baixos de discriminação.
Constata -se, assim, que a prova – no que se refere ao componente de Conhecimento Específico de Licenciatura – possuía capacidade boa de discriminar entre aqueles que dominam ou não o conteúdo.
Dentre as questões que alcançaram os maiores índices de discriminação, as de números 11 e 27 foram classificadas na categoria fácil, quanto ao índice de facilidade.
Dentre as outras nove questões com índice de discriminação muito bom, uma ( 18 ) ficou classificada como difícil e as outras oito ( 16, 17, 21, 23, 24, 25, 31 e 35 ) como de dificuldade média.
A questão de número 22 foi a mais difícil dentre as 27 questões específicas, com baixo índice de facilidade, apenas 18,0 \% de acertos.
Apesar disso, essa questão apresentou poder discriminatório médio, 0,22.
A única questão com índice de discriminação fraco, 0,15, foi a questão 26, com índice de facilidade também baixo, 0,22, o que, em termos percentuais, corresponde a 22,0 \% de estudantes que responderam acertadamente.
Pelo critério do índice de discriminação ponto bisserialo, a questão 26 foi eliminada do cômputo da nota final.
A título de exemplo das análises do comportamento das questões objetivas, o Gráfico 3.15 analisa a questão 27 do componente de Conhecimento Específico.
Esta questão foi respondida acertadamente pela maioria dos estudantes apresentando índice de facilidade 0,68, ou seja, 68,0 \% dos estudantes assinalaram acertadamente a opção D, correspondente ao gabarito.
Seu índice de discriminação foi igual a 0,56, classificado como muito bom e o maior valor alcançado por este índice para o conjunto de perguntas em análise.
Neste gráfico, cada uma das cinco curvas representa o percentual de respostas em determinada alternativa da questão 27, em função do número de acertos dos estudantes nesta parte da prova, antes de possíveis eliminações de questões pelo critério do ponto bisserial.
A alternativa correta B, representada no gráfico pela curva em verde, foi escolhida em maiores proporções pelos alunos com desempenho melhor nesta parte da prova.
Já as alternativas incorretas, também denominadas distratores, foram selecionadas principalmente por aqueles com notas mais baixas.
Aqueles com nota zero, na sua quase totalidade deixaram esta questão em branco ou marcaram mais de uma alternativa, comportamento considerado inválido.
A proporção de alunos que selecionou a resposta correta B aumenta gradativamente, chegando a atingir 100 \% para 22 acertos ou mais, enquanto a proporção dos que escolheram alternativas incorretas decai, a partir de cinco acertos, como função do número de acertos nesta parte da prova.
Os gráficos relativos às demais questões do componente de Conhecimento Específico constam do Anexo I. 3.3 ANÁLISE DAS QUESTÕES DISCURSIVAS 3.3.1 Componente de Formação Geral As análises dos resultados de desempenho dos estudantes de Geografia ( Bacharelado e Licenciatura ) nas duas questões discursivas relativas à Formação Geral encontram -se na Tabela 3.10 e no Gráfico 3.16.
Na tabela 3.10 observa -se que a nota média nesse conjunto de questões foi um pouco menor do que a obtida nas questões objetivas.
Os estudantes de todo o Brasil obtiveram, em Formação Geral, média 48,2 nas questões objetivas e 45,2 nas questões discursivas.
Pode -se notar, também, um aumento do desvio padrão de 24,3 nas questões objetivas do componente de Formação Geral dos alunos de todo o Brasil, para 30,7 nas questões discursivas do mesmo componente.
A maior média foi obtida na região Centro-Oeste ( 49,0 ) e a menor na região Sudeste ( 41,2 ).
A mediana de todo o Brasil, neste componente, foi 50,0, a mesma em todas as regiões, com exceção da região Sudeste onde a mediana foi mais baixa ( 45,0 ).
A nota máxima ( 100,0 ) e a nota mínima ( 0,0 ) foram as mesmas em todas as regiões do Brasil.
O Gráfico 3.16 representa a distribuição das notas nas questões discursivas no componente de Formação Geral segundo a opção de Bacharelado/Licenciatura.
As modas principais destas distribuições ocorrem no intervalo [ 0 ; 10 ], intervalo que inclui, além das notas zero, a frequência de alunos que deixaram este tipo de questão em branco.
A concentração de notas no intervalo modal é mais acentuada para os concluintes de Bacharelado, enquanto esta situação se inverte para a maioria dos demais intervalos.
O coeficiente de assimetria da distribuição das notas do Bacharelado é positivo ( 0,11 ), indicando assimetria à direita do gráfico, enquanto no caso da Licenciatura o coeficiente é negativo ( – 0,31 ), ou seja, a assimetria é à esquerda.
Comportamento semelhante se observa nas distribuições de notas das regiões Norte e Sudeste.
Já as distribuições de ambos os cursos nas demais regiões têm coeficientes de assimetria negativos, são assimétricas à esquerda, cauda mais acentuada no lado esquerdo do gráfico.
Na sequência, os resultados verificados para cada uma das questões discursivas de Formação Geral serão apresentados, estabelecendo -se relações com os conteúdos abordados em cada uma delas.
Os comentários da Banca de docentes corretores a respeito do observado na correção das respostas dos estudantes, suas impressões e conclusões serão apresentados junto à análise de cada questão.
Cumpre esclarecer que, tendo em vista que as questões discursivas de Formação Geral são padronizadas, ou seja, constam de todas as provas, os comentários da Banca são os mesmos para todas as carreiras acadêmicas, sendo direcionados a todos os estudantes que participaram do ENADE/2011.
A seguir, serão analisados os desempenhos da Área de Geografia ( Bacharelado e Licenciatura ) nas duas questões discursivas de Formação Geral do ENADE/2011, comparando os resultados obtidos com comentários para cada questão. 3.3.1.1 Análise da Questão Discursiva 1 do Componente de Formação Geral Os dados de Geografia ( Bacharelado e Licenciatura ), obtidos a partir das respostas à questão 1, encontram -se na Tabela 3.11 e no Gráfico 3.17.
Nessa questão – de melhor desempenho dentre as duas de Formação Geral – os alunos de todo Brasil tiveram média 48,7.
A maior média para a questão 1 foi obtida na região Centro-Oeste ( 51,5 ), e a menor, na região Sudeste ( 45,6 ).
Quanto à variabilidade das notas, o desvio padrão de todo o Brasil foi 38,0.
O menor desvio padrão foi obtido na região Centro-Oeste ( 35,3 ) e o maior desvio padrão foi obtido na região Sudeste ( 39,4 ).
As medianas das regiões Norte, Nordeste, Sul e Centro-Oeste e do Brasil como um todo foram iguais ( 50,0 ).
Na região Sudeste a mediana foi 40,0.
As notas máximas e mínimas da questão discursiva 1 foram as mesmas para todas as regiões do Brasil, respectivamente, 0,0 e 100,0.
O Gráfico 3.17 mostra a distribuição das notas na questão discursiva 1 do componente de Formação Geral para alunos do Bacharelado e da Licenciatura.
Observa -se que a maior frequência de notas dos concluintes do Bacharelado encontra -se na categoria “ em branco ”, enquanto para os alunos de Licenciatura a moda principal se encontra no intervalo ( 90 ; 100 ].
O coeficiente de assimetria é negativo e bem pequeno para a distribuição de notas da Licenciatura ( – 0,04 ) que, por isso, é praticamente simétrica.
Para o Bacharelado, no entanto, o coeficiente de assimetria é positivo ( 0,22 ) para as notas dos alunos do Brasil como um todo, como também por região, sendo o da região Norte o maior ( 0,47 ) e o da região Sul igual a zero, curva simétrica nesta última região.
No caso da notas dos alunos da Licenciatura os coeficientes são negativos e próximos de zero, sendo igual a zero na região Sudeste. 3.3.1.2 Comentários sobre a correção das respostas à Questão Discursiva 1 De maneira geral, a aplicação da chave de correção da questão 1 de Formação Geral não apresentou qualquer dificuldade digna de menção.
Para isso contribuíram, indubitavelmente, os ajustes feitos após a correção da amostra e a reunião entre todos os membros da Banca de docentes corretores.
As poucas dúvidas, todas pontuais, apresentadas pelos corretores, foram acompanhadas e respondidas pela coordenação e subcoordenação da correção das questões de Formação Geral, por meio da ferramenta de Gerenciamento de Dúvidas do Sistema de Correção On-line.
Não houve registro de qualquer ocorrência que pusesse em xeque o padrão de resposta ou a efetividade e a adequação da chave de correção.
Explica -se : trata -se de questão com comando claro, direto e objetivo ( solicitava -se, basicamente, três vantagens justificadas de cursos a distância ), cujas respostas foram corrigidas por meio da aplicação de um chave de correção testada e aprovada previamente.
Havia absoluta clareza quanto aos critérios de avaliação da correspondência entre as respostas dos estudantes e as possibilidades de vantagens de cursos a distância admitidas como corretas no padrão de resposta oficial, além de gradações explícitas ( e fáceis de aplicar ) dos diferentes níveis de pontuação previstos.
Felizmente, portanto, não há reparo a registrar em relação à facilidade de aplicação do padrão de resposta e da chave de correção, e nem em relação à atribuição dos diferentes níveis de pontuação previstos.
Todas as dificuldades que poderiam ter obstado a correta aplicação do padrão de resposta oficial e da respectiva chave de correção foram evitadas por meio dos ajustes feitos após a correção da amostra e fartamente debatidos com toda a Banca.
Digno de nota é que quantidade tão significativa de profissionais envolvidos na mesma tarefa – tanto para a questão 1 quanto para a 2 – tenha apresentado tão poucas dificuldades na execução da correção, em termos proporcionais.
Em suma, a correção da questão 1 da prova de Formação Geral do ENADE 2011 foi exemplarmente bem planejada, servindo -se de padrão de resposta muito bem adequado à questão proposta.
Quanto ao tema desta questão, em particular, a Banca verificou que uma parcela significativa de estudantes evocou experiências bastante concretas e próximas de sua realidade.
Houve várias respostas que indicavam uma vivência pessoal de ensino superior na modalidade Educação a Distância ( EaD ), evidenciando o tom de depoentes nos textos apresentados.
Os estudantes, em sua grande maioria, utilizaram parte considerável do espaço de 15 linhas disponíveis para a resposta – e outra parcela menos significativa dos que não o fizeram demonstraram notável capacidade de atender ao comando da questão de maneira objetiva, curta e, via de regra, correta.
Registre -se, ainda, que foram relativamente poucos os casos de respostas que tenham passado ao largo do tema em pauta na questão.
As capacidades de leitura, de compreensão do comando proposto e de expressão escrita que os estudantes avaliados na edição 2011 do ENADE foram satisfatórias.
Não obstante, seria leviano perder de vista que a qualidade dos textos redigidos em resposta às questões discursivas do Exame ainda está muito aquém do que se espera de concluintes de cursos de ensino superior de todas as regiões do país.
Quanto ao conteúdo das respostas, a Banca constatou boa capacidade, por parte da maioria dos estudantes, de compreensão do tema e do comando da questão.
Foram relativamente poucos os casos de respostas que deixaram de enumerar vantagens da modalidade EaD, e proporcionalmente escassos os estudantes que citaram vantagens não previstas no padrão de resposta.
Foi frequente, a tentativa direta de atender ao comando da questão.
Os erros mais comuns, em relação ao padrão de respostas e à grade de correção, foram fruto do desdobramento em vários “ itens ” daquilo que, de acordo com o padrão de respostas oficial, representava uma única vantagem.
Destacaram -se, neste caso, as respostas que apontavam a flexibilidade de horário e/ou local como duas vantagens distintas daquela modalidade de ensino.
Quanto aos diferentes níveis de pontuação previstos, a maior causa de baixas pontuações foi a ausência de justificativas, e mesmo de argumentação, para uma ou mais das vantagens enumeradas.
Isso demonstra não apenas a objetividade das respostas, coerente com a objetividade do comando da questão ( “ enumere três vantagens de um curso a distância ” ), mas também certa dificuldade de formulação plena de um texto, ou ao menos de parágrafos, em formato dissertativo – mesmo diante de uma média de 5 linhas disponíveis para cada vantagem a enumerar.
Os acertos mais comuns, ou seja, os “ itens ” do padrão de respostas mais frequentemente mencionados foram : ( 1 ) a flexibilidade de horário e/ou local ; ( 2 ) a capilaridade do ensino a distância ; ( 3 ) a democratização do acesso à educação de qualidade ; e ( 4 ) os custos menores que os de cursos presenciais.
Dentre as vantagens previstas no padrão de respostas que foram menos citadas, destacam -se a inclusão de pessoas com comprometimento motor, a qualificação de professores e a troca de experiências entre os participantes.
Foram muito comuns, no entanto, as menções à supostamente maior facilidade de acesso a professores ou/e tutores em cursos superiores a distância. 3.3.1.3 Análise da Questão Discursiva 2 do Componente de Formação Geral A Tabela 3.12 mostra que o desempenho dos estudantes na questão 2 ( média 41,4 ) foi um pouco inferior ao obtido na questão de número 1 ( média 48,7 ).
A região Centro-Oeste foi aquela onde a média, nessa questão, foi maior ( 46,4 ), e a de menor média foi a região Sudeste ( 36,8 ).
Quanto à variabilidade das notas, o desvio padrão de todo o Brasil foi 32,1, inferior ao obtido na questão de número 1 ( 38,0 ).
O maior desvio nessa questão foi obtido na região Sudeste ( 32,9 ), enquanto o menor foi obtido na região Centro-Oeste ( 29,6 ).
A mediana em todas as regiões foi a mesma que na questão 1 : diferente de 50,0 apenas na região Sudeste, onde foi 40,0.
A nota ( 100,0 ) foi obtida em todas as regiões do Brasil.
A nota mínima foi zero também em todas as regiões.
O Gráfico 3.18 mostra a distribuição das notas na questão discursiva 2 do componente de Formação Geral para alunos do Bacharelado e da Licenciatura.
Observa -se que a maior frequência corresponde aos alunos que deixaram a questão em branco tanto para aqueles de Bacharelado quanto para os de Licenciatura.
Nota -se, ainda, que a frequência de notas dos alunos de Bacharelado nas categorias “ em branco ” e “ zero ” supera a das notas de Licenciatura, enquanto para os demais intervalos a frequência de notas da Licenciatura supera ou é igual a do Bacharelado.
Existe uma moda secundária, comum às duas habilitações, no intervalo ( 50 ; 60 ].
Os coeficientes de assimetria das distribuições da Licenciatura são negativos, tanto das notas de todo o Brasil ( – 0,21 ), quanto para todas as regiões.
Para a distribuição de notas do Bacharelado o coeficiente de assimetria é positivo para o Brasil como um todo ( 0,28 ) e para as regiões Norte, Sudeste e Sul.
Nas regiões Nordeste e Centro-Oeste a assimetria da distribuição de notas do Bacharelado é negativa.
Nesta situação, os coeficientes positivos sinalizam uma semelhança à distribuição do Bacharelado ( maior concentração no zero ) e os coeficientes negativos, uma semelhança à distribuição da Licenciatura ( menor concentração no zero ). 3.3.1.4 Comentários sobre a correção das respostas à Questão Discursiva 2 Diferentemente da questão 1, a questão 2, cujo tema aborda políticas públicas para a erradicação do analfabetismo e para o aumento da empregabilidade, permitiu mais amplitude na elaboração das respostas dos estudantes, uma vez que solicita a apresentação de uma proposta.
À falta de objetividade técnica na elaboração da questão – sobretudo no seu comando – correspondeu uma miríade de possibilidades interpretativas.
Registre -se que não raras foram as respostas que se utilizaram de lugares comuns e exortações religiosas/humanitárias/cívicas, contudo, na grande maioria dos casos, essas exortações foram usadas como um encerramento do texto e não comprometeram a resposta.
Muito frequente também foi a confusão entre política pública e política partidária, bem como alguns poucos, confundiam programas educacionais com programas educativos nas TVs.
Assim, destacaram -se como propostas/programas mais recorrentes : • prosseguimento das políticas já existentes, mas com o aumento dos investimentos, normalmente sugerindo a ampliação da bolsa família, da Educação de Jovens e Adultos ( EJA ) e da EaD ; • valorização do magistério, construção de escolas e melhoria das atuais ; • parceria do governo com empresas para manutenção de salas de aula visando a alfabetização de seus funcionários, oferecimento de estágios e redução da carga horária em troca de isenção de impostos ; • parcerias com igrejas e ONGs para criação de espaços de alfabetização ; • escolas itinerantes e alfabetizadores em domicílio, principalmente para pessoas com dificuldade de locomoção, como os idosos e deficientes físicos ; • erradicação do trabalho infantil ; • vinculação da bolsa família não apenas à frequência, mas também e, principalmente, ao resultado obtido pelo aluno na escola ; • revisão das políticas atuais, sendo a mais frequente a extinção da aprovação automática ; • críticas consistentes em relação a modelos didáticos considerados inadequados e desestimulantes para a educação não só de adultos, mas de pessoas de todas as idades.
Destacam -se, ainda, outras sugestões apresentadas : • Creches nas escolas onde os pais estão sendo alfabetizados. • Diminuição de duas horas na jornada de trabalho em empresas para funcionários não alfabetizados, para que possam frequentar a escola. • Campanhas educativas vinculadas aos meios de comunicação. • Aumento do número de escolas noturnas. • Formação específica para professores alfabetizadores.
Quanto à relação entre o analfabetismo e a empregabilidade, deve -se sublinhar que nem todos os estudantes estabeleceram claramente o vínculo entre essas duas situações sociais.
Alguns falaram separadamente de uma e de outra.
Mas a maioria fez referência à necessidade de estudo para “ conseguir um bom emprego com um bom salário ”.
Alguns, em menor número, estabeleceram de forma bastante interessante a questão histórica para a situação do Nordeste ; e, também, a relação entre a pessoa analfabeta/com pouca instrução e a desempregada/empregada em serviços mais pesados e pior remunerados, que não teriam condições de educar seus filhos que, por sua vez, também não teriam melhores oportunidades no mercado de trabalho, identificando a formação de um círculo vicioso e a necessidade de sua interrupção.
Muitos afirmaram que o analfabetismo não é o único responsável pelo desemprego e sim a má distribuição de renda.
Grande também foi o número de alunos que criticaram o resultado da pesquisa, afirmando que a mesma é enganosa, uma vez que considera alfabetizada a pessoa que “ desenha seu nome ”, sendo comum associarem a atual política de alfabetização com ganhos eleitorais.
Foi comum a resposta incompleta, em que o estudante fez apenas a análise das desigualdades/crítica do quadro apresentado, ou só apresentou proposta.
Alguns se limitaram a responsabilizar o governo referindo -se de forma bastante genérica à questão da “ educação ” e “ profissional ”.
Entre os equívocos que mais se repetiram, destacou -se a simples análise dos dados apresentados na tabela que consta do enunciado.
Alguns textos, inclusive, estavam corretos, sem que, no entanto, fosse respondida a pergunta.
Da mesma forma, quando os estudantes partiram para a segunda parte da questão, fizeram referência a vários programas já existentes ou simplesmente disseram que é muito importante que existam projetos para a educação.
Outro equívoco recorrente foi a análise da educação no Brasil como um todo.
Essa análise, apesar de correta, não se referia especificamente à questão do analfabetismo, o que, por conseguinte, levava a sugestões que não eram direcionadas à erradicação do analfabetismo, e sim à melhora da educação no Brasil.
Sendo assim, foram apresentadas propostas como o aumento do número de faculdades, o reforço de alunos do Ensino Médio, a criação de escolas técnicas, etc.
Por outro lado, houve um grande número de redações bem escritas e precisas, no que se refere ao que foi exigido pela questão.
Foram análises equilibradas e sensatas, correlacionando a problemática do analfabetismo com o desemprego e a apresentação de sugestões bastante consistentes. 3.3.2 Componente de Conhecimento Específico Na parte da prova relativa às questões discursivas no componente de Conhecimento Específico ( Tabela 3.13 ), observa -se que a média foi mais baixa do que para as questões discursivas do componente de Formação Geral.
Enquanto no componente de Formação Geral a média para estudantes de Geografia de todo o Brasil foi 45,2, na parte de Conhecimento Específico a média foi 11,1.
A maior média deste componente foi obtida pelos estudantes da região Sul ( 11,8 ), e a menor, pelos da região Norte ( 9,6 ).
Quanto à variabilidade das notas, o desvio padrão de todo o Brasil foi 14,4.
O maior desvio padrão foi encontrado na região Sudeste ( 15,5 ), e o menor, na região Norte ( 11,9 ).
A maior nota máxima foi obtida nas regiões Nordeste, Sudeste e Sul ( 86,7 ) enquanto a menor nota máxima foi encontrada na região Centro-Oeste ( 73,3 ).
Além disso, a nota mínima ( 0,0 ) foi obtida por alunos de todas as regiões do Brasil sem exceção.
A mediana do Brasil como um todo foi 5,0, também encontrada na região Norte.
As demais medianas foram : 3,3 nas regiões Nordeste e Sudeste ; e 6,7 nas regiões Sul e Centro-Oeste.
O Gráfico 3.19 representa a distribuição das notas nas questões discursivas no componente de Conhecimento Específico segundo a opção de Bacharelado/Licenciatura.
As modas destas distribuições ocorrem no intervalo [ 0 ; 10 ].
O percentual de alunos de Bacharelado com notas no intervalo [ 0 ; 10 ] é maior do que o de Licenciatura, enquanto esta proporção se inverte nos dois intervalos seguintes.
Nota -se que ambas as distribuições são decrescentes até o intervalo ( 80 ; 90 ], onde já ocorrem pouquíssimas notas.
Este comportamento das distribuições, assimetria positiva acentuada, é confirmado pelos coeficientes de assimetria positivos e maiores do que 1 para os dois cursos e em todas as regiões.
No Gráfico 3.19 os coeficientes de assimetria são 1,54 e 1,53 para Bacharelado e Licenciatura, respectivamente.
A análise de cada uma destas questões será feita a seguir. 3.3.2.1 Análise da Questão Discursiva 3 do Componente de Conhecimento Específico Na questão 3, cujos resultados aferidos encontram -se descritos na Tabela 3.14, a média dos estudantes de todo o Brasil foi 17,3.
A menor média nessa questão foi obtida pelos alunos da região Norte ( 15,7 ), enquanto a maior média foi obtida na região Centro-Oeste ( 18,5 ).
Quanto à variabilidade das notas, o desvio padrão de todo o Brasil foi 22,6.
O maior desvio padrão foi obtido na região Sul ( 23,4 ), enquanto o menor foi obtido na região Norte ( 20,9 ).
A nota máxima, 100,0 pontos, foi alcançada por pelo menos um aluno da região Sudeste.
Nas demais regiões a nota máxima foi 90,0.
A nota mínima ( 0,0 ) foi obtida em todas as regiões do Brasil, sem exceção.
A mediana do Brasil como um todo foi 0,0, o que indica que mais do que 50 \% dos alunos obtiveram nota zero ou não resolveram essa questão.
A única região com mediana diferente de zero foi a Centro-oeste, onde esta estatística foi 10,0.
O Gráfico 3.20 mostra a distribuição das notas na questão discursiva 3, do componente de Conhecimento Específico da área de Geografia para alunos do Bacharelado e da Licenciatura.
Essa distribuição tem moda nas questões em branco, com percentual maior para os alunos de Bacharelado.
O coeficiente de assimetria é positivo e em torno de 1 para os dois cursos : 1,11 para o Bacharelado e 1,00 para a Licenciatura.
A assimetria é positiva para todas as distribuições por região, sendo que a distribuição das notas dos alunos de Bacharelado da região Norte tem o maior coeficiente ( 1,60 ) e o do curso de Licenciatura da região Centro-Oeste o menor ( 0,82 ).
Esta proximidade nos valores dos coeficientes de assimetria sugere que a forma dos histogramas é bem semelhante para todas as combinações de região e habilitação. 3.3.2.2 Comentários sobre a correção das respostas à Questão Discursiva 3 A questão trata da crise atual europeia e suas repercussões no contexto contemporâneo buscando aferir, a partir da apresentação de uma charge, as competências e habilidades relacionadas a : leitura, interpretação, análise crítica, capacidade argumentativa, elaboração de sínteses e de avaliação prospectiva sobre um tema da atualidade.
Assim, no item “ a ” é solicitado que o estudante faça uma abordagem geral sobre a crise socioeconômica européia e suas repercussões junto à população.
No item “ b ” solicita -se que trate das relações entre os países da Europa representados na charge ( Alemanha, França e outros países que estão em situação pior diante da crise ).
No item “ c ”, solicita -se que discorra sobre aspectos da relação entre os países europeus e o crescimento da economia dos países emergentes.
Não resta dúvida que a questão é atual e abrangente e deixou muito claro o que pretendia avaliar.
A única \textbf{observação} crítica é que os acontecimentos sobre a temática estão em curso e ainda não se tem clareza de todos os desdobramentos e aspectos específicos relacionados ao tema.
Por exemplo, o crescimento das economias dos países emergentes, nesse contexto de crise no centro da dinâmica capitalista, ainda não é claramente diagnosticado.
Daí surgiu a necessidade de abrir a resposta do item “ c ” para muitas possibilidades.
A questão foi bem elaborada e parte do princípio de que as temáticas espaciais da geografia regional devem ser trabalhadas na formação básica do futuro geógrafo/professor de Geografia.
Tal perspectiva vem ao encontro da necessidade de que os futuros profissionais da Geografia tenham uma leitura competente sobre os modelos de regionalização, sobre a estrutura de poder dos Estados nacionais europeus, sobre as crises cíclicas do capitalismo internacional e sobre os impactos sociais nos países envolvidos na referida crise mundial e regional.
A proposta deve ter levado em conta a necessidade de compreensão, por parte de futuros geógrafos e professores de Geografia, sobre a ação de potências regionais e de países semiperiféricos, no sistema-mundo, bem como sobre as estruturas e concepções que devem fazer parte da compreensão espacial dos Estados nacionais contemporâneos.
Para a questão 3 ser bem respondida, o estudante deveria demonstrar uma boa formação geral, tendo realizado leituras consistentes ao longo da sua formação acadêmica.
Também era uma necessidade, para uma boa resposta nessa questão, o acompanhamento do que é informado pelos meios de comunicação, no sentido do compromisso de observar as transformações conjunturais e estruturais do mundo contemporâneo, de forma a ter possibilidade de articular as mudanças socioespaciais às bases teóricas da Geografia regional, política e econômica.
Os professores que participaram da correção consideraram que a aplicação do padrão de resposta tornou a correção relativamente ágil e cômoda.
No entanto, especialmente nos itens “ b ” e “ c ”, houve necessidade de uma abertura maior de aspectos a considerar.
Apesar disso, vários dos estudantes não conseguiram chegar às expectativas apresentadas como padrão de respostas.
Lacunas e inadequações pareceram se manifestar, mais por falta de conteúdo e dificuldade na elaboração de texto escrito, do que pela falta de compreensão em relação ao que se desejava aferir.
O desempenho dos estudantes nessa questão mostrou -se fraco, apesar de grande parte dos concluintes que apresentaram alguma resposta para a questão construírem respostas que mostravam o entendimento do que estava sendo solicitado.
O tipo de resposta mais frequente foi aquele em que o estudante fez um comentário sobre a crise e apontou o desemprego como efeito para as populações.
Das respostas obtidas pôde -se perceber e concluir que os estudantes : 1 ) não acompanham os jornais diários ( impressos e na televisão ) e nem se atualizam sobre os temas da conjuntura internacional na internet ; 2 ) parecem ter pouco contato com as temáticas regionais da Geografia nos cursos de formação geral ; 3 ) demonstraram conhecer pouco o tema “ regionalização do mundo ” e, portanto, uma das temáticas estruturantes da Geografia ( a “ região ” ) está sendo deixada de lado na formação do geógrafo / professor de Geografia ; 4 ) confundem países emergentes da periferia global ( Brasil, China, Rússia... ) com os países da periferia comunitária europeia ( Irlanda, Grécia, os países do leste europeu... ) ; 5 ) que apresentaram as melhores abordagens regionais o fizeram por meio de textos marcados por um discurso frequente nos anos 70 e 80 da Geografia moderna : a tensão entre o capitalismo e o socialismo e o sistema de submissão do Estado “ às forças do capital ”, com um viés marxista ortodoxo muito forte e pouco esclarecedor ; as análises que envolvem, além da dimensão econômica, as dimensões do simbólico e do cultural foram muito escassas. 3.3.2.3 Análise da Questão Discursiva 4 do Componente de Conhecimento Específico A Tabela 3.15 contém as informações relativas à questão 4 do conjunto de questões do componente de Conhecimento Específico.
O desempenho dos estudantes de todo o Brasil nesta questão foi inferior ao desempenho na questão de número 3.
A média geral do Brasil foi 7,2, sendo a menor média registrada na região Norte ( 5,8 ) e a maior na região Sul ( 8,1 ).
A nota máxima ( 100,0 ) foi atingida em todas as regiões.
A mediana e a nota mínima foram 0,0 em todas as regiões do Brasil, sem exceção.
Este fato indica que mais de metade dos estudantes de Geografia em toas as Grandes Rgiões tirou nota zero nesta questão.
O Gráfico 3.21, representa a distribuição de notas na questão discursiva 4, no componente de Conhecimento Específico para alunos do Bacharelado e da Licenciatura.
A maior frequência, tanto para estudantes de Bacharelado quanto para Licenciatura, ocorre na categoria “ em branco ”, seguida de perto pela nota zero.
Observa -se que uma maior proporção de estudantes deixou essa questão em branco, comparando -se à questão discursiva 3, como também o percentual de notas zero é bem superior nesta questão.
Os coeficientes de assimetria são positivos e bastante altos : 3,08 para a distribuição de notas dos concluintes de Bacharelado e 2,75 para os de Licenciatura.
Este comportamento se repete para ambos os cursos em todas as regiões, sendo o menor coeficiente de assimetria igual a 2,38, da distribuição de notas de Licenciatura da região Sul, e o maior, 4,44, do Bacharelado da região Norte.
Novamente cumpre notar que esta proximidade nos valores dos coeficientes de assimetria sugere que a forma dos histogramas é bem semelhante para todas as combinações de região e habilitação. 3.3.2.4 Comentários sobre a correção das respostas à Questão Discursiva 4 A Banca de docentes corretores considera, após análise da questão e do processo de correção, que a questão 4 foi elaborada de forma confusa.
Os problemas relacionados às dificuldades para responder podem estar na pouca objetividade do comando da questão, muito extenso.
Caso tivesse sido bem elaborada, seria possível considerá -la uma boa maneira de verificar as competências e habilidades que um estudante de Geografia deve ter.
Trata -se de aferir a capacidade de estabelecer relação entre conceitos e conteúdos de Geografia, sobre a relação entre natureza e sociedade, além de aferir competências e habilidades de aplicação dessas relações a partir da elaboração de um plano de aula.
A questão 4 solicita aos graduandos a elaboração de um texto expondo a importância da promoção da articulação, no ensino escolar da Geografia, de uma questão de caráter epistemológico crucial no âmbito da ciência geográfica e que, indubitavelmente, distingue -a de outros ramos do conhecimento científico : a articulação entre os fenômenos físicos ( clima, relevo, vegetação, hidrografia e solos ) e os fenômenos humanos ( econômicos, políticos e sociais ).
Trata -se da chamada integração sociedade-natureza inerente aos estudos geográficos. sociais ).
Trata -se da chamada integração sociedade-natureza inerente aos estudos geográficos.
Ademais, nessa mesma questão, outro aspecto de cunho epistemológico essencial ao conhecimento geográfico é exigido : a articulação dos conteúdos da Geografia ( estudos sobre indústria, população, agricultura, urbanização, geomorfologia, pedologia, climatologia, cartografia e outros ) aos conceitos que legitimam e, por consequência, garantem a cientificidade dessa disciplina ( espaço, território, região, paisagem e lugar ).
Estas categorias geográficas, quando conceituadas, embasam a Geografia e lhe dão a credibilidade necessária ao seu fazer científico.
Houve a exigência, também, de elaboração de um plano de aula esquemático, o qual deveria conter os tópicos indispensáveis a um planejamento que envolvesse a relação “ conteúdo – conceito ”.
Para muitos estudantes, essa questão se referia a um tema específico para os alunos da licenciatura, o que constituiu um equívoco, visto que a prova foi aplicada para todos os estudantes concluintes do curso de Geografia, dentre os quais muitos cursaram apenas o bacharelado.
Daí, comentários do tipo : “ No bacharelado não se aprende a fazer plano de aula ”.
Na avaliação da Banca, não foi uma escolha adequada o fragmento de um artigo que foi utilizado na questão, uma vez que foi dissociado do contexto geral do artigo.
O trecho, originalmente, está na introdução do artigo da referência, e tem como objetivo levantar os vários aspectos que são abordados ao longo do trabalho.
Em função disso, na condição de suporte da questão, o trecho apresentava um leque de temas relacionados ao ensino da Geografia.
Some -se a isso o fato de que o comando da questão propõe algo inviável de ser contemplado, até mesmo se a proposta fosse direcionada a professores com muita experiência.
Construir um texto com as solicitações feitas, abordando todos os aspectos indicados, não era plausível para o tempo da prova e espaço destinado à resposta dessa questão.
Diante da análise aqui exposta, a Banca adotou o seguinte procedimento para a avaliação : considerar uma das duas possibilidades de resposta, isto é, ou a articulação dos conteúdos aos conceitos enfocando a relação entre os aspectos de ordem física e humana, por meio de um texto, ou esse enfoque se restringindo à apresentação de um plano de aula esquemático.
Feito isso, a resposta decomposta numa das possibilidades seria avaliada de zero ( 0 ) a dez ( 10,0 ).
O desempenho dos estudantes foi ruim.
Muitos discorreram sobre a importância da Geografia, sobre a superação de uma Geografia tradicional por uma Geografia crítica, sobre como deve ser o ensino de Geografia, ou como deve ser uma aula etc., sem tocar no que estava sendo requerido.
Além disso, muitos estudantes escreveram sobre a importância do trabalho de campo, fizeram críticas aos livros didáticos ou tentaram explicar como uma aula poderia ser melhor.
Sobre este último tópico, chamou a atenção o grande número de respostas que associava um bom ensino às práticas que resultam em"uma aula alegre","lúdica","divertida", com jogos e brincadeiras, sendo que, na maioria das vezes, não fizeram referência alguma à Geografia, mas sim a uma aula genérica.
Um tipo de resposta recorrente, que também não respondeu ao solicitado, os estudantes apenas definiram alguns conceitos da Geografia.
Outros, sem o domínio dos principais conceitos da disciplina, elencaram como conceitos da Geografia, por exemplo,"cartografia","vegetação","relevo","hidrografia","miséria", “ climas do Brasil","a Terra","Rio São Francisco"etc.
Outros, ainda, montaram um esquema de plano de aula sem tocar na relação entre os conceitos e os conteúdos, sendo que foi frequente a montagem de um plano de aula sem qualquer vínculo com o que foi requerido na questão ou com temas da Geografia, apenas apresentando as etapas da aula.
Alguns concluintes tentaram explicar a relação entre os conceitos e conteúdos, mas acabaram escrevendo sobre a relação em si, sem explicitar qual o conceito.
Um ponto positivo foi o fato de que um número expressivo de estudantes, apesar de não ter entendido o enunciado, definiu de forma bastante precisa os principais conceitos da Geografia, o que reflete a preocupação de alguns cursos em trabalhar a questão conceitual da disciplina.
Das respostas corrigidas, o que ficou evidente foi a ausência do estabelecimento de conexões entre os aspectos físicos e humanos e entre os conteúdos e os conceitos geográficos.
Muitos, além de não estabelecerem conexões, mostraram -se incapazes de realizar, ainda que de modo precário, inferências quanto a isso na resposta.
Como resultado, o quantitativo de \textbf{repostas} fracas foi muito elevado, das quais várias com nota zero.
Desconsiderando -se o número elevadíssimo de respostas em branco, dentre as corrigidas, foi possível constatar respostas que, além de desconexas, apresentavam erros graves de redação.
Muitos dos candidatos ao exercício do magistério em Geografia, em razão dos erros ortográficos gravíssimos constatados ( alguns possivelmente por reproduzirem a linguagem oral ), não demonstraram ter condições minimamente aceitáveis para uma prática docente satisfatória.
Já as respostas com alta pontuação compuseram um universo bastante restrito.
Houve um contingente muito pequeno de estudantes que conseguiu realizar as conexões exigidas na questão 4 ( entre os aspectos físicos e humanos e entre conceito e conteúdo ) e demonstrou domínio sobre o que foi solicitado.
Nesses poucos casos, os textos também apresentavam inteligibilidade na exposição do argumento. 3.3.2.5 Análise da Questão Discursiva 5 do Componente de Conhecimento Específico A Tabela 3.16 contém as informações relativas à questão 5 do conjunto do componente de Conhecimento Específico.
O desempenho dos estudantes nessa questão foi tão fraco quanto nas questões 3 e 4.
A nota média dos estudantes de todo o Brasil foi 8,6.
A maior média foi registrada na região Sudeste ( 9,8 ), enquanto a menor média foi registrada na região Norte ( 7,1 ).
Quanto à variabilidade das notas, o desvio padrão dos alunos do Brasil, como um todo, foi 18,6.
Enquanto o maior desvio foi encontrado na região Sudeste ( 20,5 ), o menor foi encontrado na região Norte ( 15,4 ).
A nota máxima ( 100,0 ) foi alcançada em todas as regiões.
A mediana e a nota mínima foram zero para todas as regiões do Brasil, ou seja, pelo menos metade dos alunos tirou nota zero na questão.
O Gráfico 3.22 representa a distribuição das notas da questão discursiva 5 no componente de Conhecimento Específico para alunos do Bacharelado e da Licenciatura.
Mais uma vez, destaca -se o grande número de estudantes que deixaram a questão em branco, para duas habilitações desta Área, mais do que 50 \%, como já indicava o valor da mediana.
Como nas questões discursivas anteriores do Componente de Conhecimento Específico, esta tem assimetria positiva grande, neste caso igual ou maior do que 2, o que explica a longa cauda à direita do gráfico.
Para alunos do Bacharelado o coeficiente de assimetria é de 2,20 e para Licenciatura de 2,79.
Este tipo de assimetria também ocorre em todas as regiões, sendo o maior coeficiente 3,04, da distribuição de notas da licenciatura da região Nordeste. 3.3.2.6 Comentários sobre as respostas à Questão Discursiva 5 Essa foi considerada pela Banca, sem dúvida, a questão mais difícil da prova.
O comando da questão é pouco claro.
Por exemplo, quando há referências no comando às atuais correntes de pensamento geográfico, surge a dúvida : o que é atual para o autor da questão?
O atual está se referindo às últimas décadas ou é uma referência mais recente, dos últimos anos?
O autor da questão está se referindo ao momento de pensamento desde a Geografia Quantitativa, desde a Geografia Crítica ou está se referindo à Geografia Cultural, ou à Geografia Humanística?
O padrão de resposta não fazia nenhuma menção à relação entre as correntes de pensamento geográfico “ atuais ” e o geoprocessamento, o que fazia parte do que foi perguntado.
O texto da questão, por conseguinte, estava confuso e o padrão de resposta inicialmente apresentado não contemplava a proposta.
Apesar da pertinência da temática, o comando da questão é demasiado abrangente, o que pode ter interferido na qualidade das respostas, especialmente considerando o tempo e o espaço destinado às mesmas.
Desta forma, faltaram objetividade e delimitação clara dos pontos a serem abordados, que ora focava apenas o Geoprocessamento, ora apenas o saber geográfico, suas correntes e conceitos.
Cabe ressaltar, no entanto, o mérito da temática.
A proposição de uma questão sobre a área de Geoprocessamento é importante em função do estágio de desenvolvimento técnico da sociedade e de grande relevância para a ciência geográfica na atualidade.
O Geoprocessamento, como área de conhecimento, processa diversos tipos de informações geográficas por meio de técnicas matemáticas e computacionais.
O Sistema de Informação Geográfica ( SIG ), como uma ferramenta do Geoprocessamento, permite gerar informações pela análise e integração de dados geográficos, criando diversos mapas temáticos.
Estes facilitam a análise, gestão e representação do espaço geográfico, subsidiando tomadas de decisão.
O Geoprocessamento tem influenciado fortemente as áreas de Cartografia, Análise de Recursos Naturais, Transportes, Comunicações, Energia e Planejamento Urbano e Regional.
Assim, para os profissionais da Geografia, na atualidade, é fundamental o conhecimento e formação na área de Geoprocessamento, e é interessante provocar a reflexão sobre esta temática.
As respostas apresentadas apontaram o estágio inicial da utilização e do domínio do Geoprocessamento no sistema universitário da Geografia.
Apesar da atualidade e da importância da temática, foi observado que a utilização e o domínio do Geoprocessamento ainda são muito incipientes.
Dentre as respostas apresentadas os alunos buscaram avaliar as instituições, uma vez que apontavam para o conhecimento da estrutura curricular das universidades e sua capacidade de infraestrutura técnica e mão de obra especializada para o Geoprocessamento.
Na realidade, o Geoprocessamento é uma área de conhecimento e o SIG, um sistema/ferramenta do Geoprocessamento.
Poucas respostas mencionaram que o conhecimento é produzido a partir das técnicas e da Matemática, assim como a importância do Geoprocessamento em outras áreas.
Quase nunca foi mencionada a relação do Geoprocessamento e do SIG com pensamento geográfico.
Foram verificados muitos problemas na redação e várias respostas em branco.
Poucas respostas boas e claras foram apresentadas, predominando respostas fracas e algumas medianas.
Textos com a menção à importância do Geoprocessamento e à falta de formação e qualificação dos professores na área foram bastante comuns.
Dessa forma, os estudantes questionaram a qualidade da formação que receberam e revelaram que há poucos docentes que atendam aos novos campos do conhecimento que dão suporte aos estudos geográficos.
As respostas mencionaram, também, a dificuldade de acesso aos programas e equipamentos, assim como a complexidade da linguagem dos programas.
Foram comuns respostas que apontaram para a falta de infraestrutura técnica e para a dificuldade de acesso à tecnologia, tanto nas escolas quanto nas universidades.
Houve uma parcela significativa de respostas que demonstraram que os estudantes apenas ouviram falar no Geoprocessamento, assim como sua importância na atualidade.
No entanto, foram \textbf{repostas} curtas, que não apresentaram ou desenvolveram nenhuma outra informação, revelando o pouco conhecimento que possuem sobre o que realmente significa o Geoprocessamento, suas possibilidades de uso e restrições.
Muitos concluintes também revelaram o pouco conhecimento sobre as correntes geográficas, respondendo de forma errada ou mesmo evitando seguir esse caminho de resposta, optando por tratar apenas do Geoprocessamento ou vinculá -lo ao ensino em sala de aula. 3.3.3 Considerações Finais da Banca sobre a correção das questões discursivas do Componente de Conhecimento Específico Buscou -se apresentar uma síntese das principais \textbf{observações} feitas no decorrer da correção.
Algumas, no entanto, se referem a características que valem a pena ser ressaltadas para a avaliação do curso de Geografia no ENADE/2011. 1 ) As temáticas das questões são bastante relevantes para os estudos espaciais, embora nem todas as questões tenham sido bem elaboradas. 2 ) Os estudantes ( excetuando -se cerca de 20 \% deles ) demonstram que não levam a sério a elaboração de respostas às questões do ENADE, escrevendo sem nenhum cuidado, com a letra ilegível, o que dificulta a correção.
Além disso, respondem de maneira muito superficial e generalista a questões que mereceriam uma análise mais profunda e complexa. 3 ) Há deficiência na habilidade de construção do texto escrito para além da correção gramatical ; as dificuldades de articulação de ideias que envolvem teoria / conceito / evento são muito presentes, o que denota grande deficiência na formação acadêmica. 4 ) Sem dúvida, não se pode deixar de considerar o \textbf{peso} das falhas nos enunciados das questões 4 e 5.
É possível que o elevado número de respostas que não conseguiram empreender, por exemplo, a relação proposta pela questão 4, entre os aspectos físicos e humanos e entre os conteúdos e conceitos, tenha como origem a falta de clareza necessária do enunciado. 5 ) Outro aspecto a considerar : a dificuldade nas questões 4 e 5 manifestadas nas reclamações dos graduandos em relação ao fato de que elas não se enquadravam ao perfil deles ( a questão 4 não poderia ser endereçada aos concluintes do bacharelado e a questão 5 não poderia ser endereçada aos concluintes de licenciatura ) ; talvez esse seja um motivo que contribuiu para respostas equivocadas.
Na questão 4, é preciso acrescentar o fato da proposta de cobrança versar sobre uma empreitada pouco exequível para o espaço destinado à resposta.
Ainda que os corretores estivessem considerando ou o texto ou o plano de aula, enquanto modalidades separadas da avaliação, isso não era de conhecimento dos que responderam a questão.
Portanto, presume -se que dúvidas e dificuldades tenham sido geradas em razão da formulação da questão.
Obviamente, essa consideração em nada isenta aqueles que, pelas respostas dadas, demonstraram profundas dificuldades de expressão escrita, denotando que as dificuldades no \textbf{aprendizado} em Geografia de muitos graduandos que compuseram este universo de avaliação advêm das deficiências na formação da Educação Básica, a mesma na qual parte deles pleiteia atuar, como professores, ao terminar o curso. 6 ) Sobre a questão 5, houve reconhecimento, por parte dos estudantes, da importância do Geoprocessamento para a Geografia e o acesso limitado a essa área de conhecimento.
De certa forma, a questão 5 serviu ao propósito de avaliar as instituições porque revelou as deficiências na estrutura curricular das universidades e na sua capacidade de infraestrutura técnica e mão de obra especializada para o Geoprocessamento. 7 ) Ficou demonstrada uma perspectiva crítica por parte dos concluintes sobre a qualidade dos cursos de Geografia, especialmente na diferenciação de qualidade na formação entre Licenciatura e Bacharelado.
Tal percepção revela o preconceito que pesa contra os profissionais do ensino, que além de terem seus cursos menos valorizados socialmente, têm menor número de disciplinas específicas.
Por fim, registra -se a dificuldade no uso formal da língua portuguesa escrita, evidenciando problemas anteriores aos cursos de graduação, não resolvidos pelas instituições de ensino superior.
Isso aponta para a premência da adoção de medidas que possam combater essa deficiência ainda nos primeiros anos da Educação Básica.
Além desse aspecto, as \textbf{observações} da Banca de corretores Geografia do ENADE/2011 servem para, de modo contundente, endossar as enormes dificuldades da grande maioria dos graduandos da Área em não apresentar domínio ( de conteúdo e de competências cognitivas básicas ) para formular respostas adequadas às questões discursivas da prova.
Nem a flexibilização ampla no padrão de resposta originalmente estabelecido ( ainda que advindo do reconhecimento das falhas na elaboração de algumas questões ) foi suficiente para minorar as respostas de baixa pontuação.

% matched lemmas: aprendizado, observação, peso, repor
\end{textsample}
